%%%%%%%% MICE: Mixture of In-Context Experts for recurring concept drift. ICML template (icml2026.sty until the
%%%%%%%% 2027 style is released). Skeleton of 2026-10-03: no experiment has been run; every number is a TODO.
\documentclass{article}

\usepackage{microtype}
\usepackage{graphicx}
\usepackage{booktabs}
\usepackage{hyperref}
\newcommand{\theHalgorithm}{\arabic{algorithm}}

\usepackage{icml2026}          % blind review
% \usepackage[accepted]{icml2026}

\usepackage{amsmath}
\usepackage{amssymb}
\usepackage{mathtools}
\usepackage{amsthm}
\usepackage[capitalize,noabbrev]{cleveref}
\usepackage{algorithm}
\usepackage{algorithmic}
\usepackage[textsize=tiny]{todonotes}

\theoremstyle{plain}
\newtheorem{theorem}{Theorem}
\newtheorem{proposition}{Proposition}
\newtheorem{lemma}{Lemma}
\theoremstyle{definition}
\newtheorem{definition}{Definition}
\newtheorem{assumption}{Assumption}
\newtheorem{observation}{Observation}
\theoremstyle{remark}
\newtheorem{remark}{Remark}

\input{math.tex}
\crefname{assumption}{Assumption}{Assumptions}
\crefname{observation}{Observation}{Observations}
\newcommand{\method}{\textsc{MICE}}   % working name: Mixture of In-Context Experts

\begin{document}

\twocolumn[
  \icmltitle{Recall or Relearn? A Mixture of In-Context Experts for Recurring Concept Drift}
  \icmlsetsymbol{equal}{*}
  \begin{icmlauthorlist}
    \icmlauthor{Anonymous Author(s)}{anon}
  \end{icmlauthorlist}
  \icmlaffiliation{anon}{Anonymous Institution}
  \icmlcorrespondingauthor{Anonymous Author(s)}{anon@anon.org}
  \icmlkeywords{concept drift, recurring concepts, tabular foundation models, in-context learning, online learning}
  \vskip 0.3in
]
\printAffiliationsAndNotice{}

\begin{abstract}
Tabular foundation models (TFMs) learn a data stream through their context, so a past concept can be recalled by
putting its rows back into the context, at no training cost. We ask when recalling a concept is worth more than
relearning it. Context-only stream learners forget recurring concepts, but a TFM also relearns a concept from a few
hundred rows, so a reset-on-drift detector already recovers most of the loss. We show that the value of memory is set
by the TFM's learning curve: the gain of recall over relearning equals the area between the learning curve and its
value at the context budget over the recurrence, minus one batch of identification. We propose \method{}, a mixture of
in-context experts that combines fast windows with stored concepts through exponential weights. On a controlled grid,
the gain predicted from learning curves measured beforehand ranks the observed gains with Spearman 0.94 on held-out
seeds, memory beats resetting on every stream with hard concepts (up to 8.8 points), and a window ensemble without
memory does not; on the standard drift generators, which a TFM relearns within 200 rows, the predicted and observed
gains are both near zero.
\todo{Add the real-stream results.}
\end{abstract}

% =====================================================================================================
\section{Introduction}
\label{sec:intro}

\textbf{Momentum.} TFMs such as TabPFN \citep{hollmann2025tabpfn} and TabICL \citep{qu2025tabicl} approximate
posterior prediction in one forward pass \citep{muller2022pfn,nagler2023statistical}. In stream mining they replace
incremental trees: a TFM with a sliding memory outperforms adaptive tree ensembles on non-stationary benchmarks
\citep{lourenco2026streams}, and a prior with drifting causal models improves TFMs under temporal shift
\citep{helli2024drift}.

\textbf{Problem.} Real streams are non-stationary \citep{gama2014survey,lu2019learning}, and many of their concepts
\emph{recur}: seasons, weekdays, market regimes, sensor modes. For a trained model a recurring concept means
retraining or keeping a pool of models; for an in-context learner it only means putting the right rows back into the
context. Existing in-context stream learners do not exploit this. A FIFO memory forgets a concept once it leaves the
window, and a drift-extrapolating prior assumes a trend rather than a return.

\begin{observation}\label{obs:forget}
FIFO contexts forget recurring concepts. In the five batches after a concept recurs, a TFM with a FIFO context or
with the dual memory of \citet{lourenco2026streams} reaches 0.16--0.52 accuracy on Sine, Stagger and Agrawal streams,
whereas the same TFM with the context of the matching past concept reaches 0.99--1.00; in steady state all contexts
are equally accurate. \todo{Dev seeds; replicate on test seeds.}
\end{observation}

\begin{observation}\label{obs:relearn}
Relearning in context is cheap for simple concepts. Resetting the context when a drift detector fires
\citep{gama2014survey} closes most of this gap: on Stagger it reaches the ceiling set by the one-batch label delay,
and no memory can improve on it. Memory pays only when a concept needs many rows to be relearned: on mixtures with
100 centroids, a TFM needs more than 1000 rows to approach its asymptote, and recalling the past context beats
resetting by 8.6 points with blocks of 500 rows and by 2.6 points with blocks of 2000 rows on held-out seeds, while the
gain is about zero for the standard drift generators, whose concepts a TFM relearns within about 200 rows.
\end{observation}

This raises the question: \emph{when should an in-context learner recall a past concept rather than relearn it, and
how can it do both without knowing which concept is active?}

\textbf{Approach.} We treat the context of every past interval as an offline expert and the context of the current
interval as an online expert, and combine them with a meta expert, following the mixture of online and offline
experts (MOOE) of \citet{zhao2025mooe}. Unlike MOOE, an in-context expert costs no training: it is a stored context.
The experts are weighted by their losses on the most recent labelled batch; fast windows
next to the stored experts cover concepts that are new or cheap to relearn.

\textbf{Contributions.}
(1) A definition of recurring concept drift for in-context learners, in which past contexts are free experts
(\cref{def:recur}).
(2) Observations that FIFO contexts forget recurring concepts, that in-context relearning is cheap for simple
concepts, and that the value of memory is set by the learning curve.
(3) A recall-versus-relearn proposition that predicts the value of memory from the learning curve, confirmed
quantitatively on held-out streams.
(4) \method{}, a training-free mixture of window and stored in-context experts, and the finding that the standard
drift generators are too easy for TFMs to evaluate memory.

% =====================================================================================================
\section{Recurring Concept Drift for In-Context Learners}
\label{sec:def}

A stream $\cS_{0,T}$ is split into offline intervals $I_1,\dots,I_K$ with distributions $P_1,\dots,P_K$ and an
online interval $I_{K+1}$ that keeps receiving rows. Concept drift occurs between intervals in the sense of
\citet{lu2019learning}.

\begin{definition}[Recurring in-context concept drift]\label{def:recur}
Let $\cC_k$ be a context sampled from interval $I_k$ and $f_{\Theta}(\cdot\,;\cC_k)$ the frozen TFM conditioned on it
(the $k$-th \emph{in-context expert}). The drift into the online interval is \emph{recurring} if
$\min_{k\le K}\Div(P_{K+1},P_k)\le\epsilon$ for a discrepancy $\Div$ and tolerance $\epsilon$, and \emph{new} otherwise.
\end{definition}

\todo[inline]{Choose $\Div$: I-Div (label-free, current window) and R-divergence (model-oriented, labelled windows)
\citep{zhao2024idiv,zhao2023rdiv}.}

% =====================================================================================================
\section{Method}
\label{sec:method}

\method{} runs a frozen TFM with batches of $B$ rows, labels one batch late and contexts of at most $M$ rows.

\textbf{Experts.} Two kinds of in-context experts predict every batch: \emph{windows}, FIFO contexts of the last 100,
300 and $M$ labelled rows (fast relearning), and \emph{stored concepts}, contexts kept in a pool (memory). Every 500
labelled rows close a segment. The segment is merged into the stored expert whose context predicts it at least as well
as the segment predicts itself by 2-fold cross-fitting, minus a tolerance of 0.02, in the spirit of the
model-oriented discrepancy of R-divergence \citep{zhao2023rdiv}; otherwise it becomes a new stored expert. A stored
expert keeps its last $M$ rows; the pool holds at most 12 experts (least recently used out).

\textbf{Meta expert.} The prediction is the mixture of the experts' predictive distributions with weights
$w_k\propto\exp(-\eta\,\ell_k)$, where $\ell_k$ is the log-loss of expert $k$ on the most recent labelled batch
(exponential weights over in-context experts, as the meta expert of MOOE \citep{zhao2025mooe}). The constants
($\eta=10$, segment length 500) were chosen on development streams.

No expert is trained: storing a concept means storing rows, and recalling it means conditioning on them.
\todo{Label-free prior weights from a two-sample discrepancy (I-Div) were planned but are not part of the evaluated
method; the evaluated streams with recurrence have identical $P(\vx)$ across concepts, where no label-free signal
exists.}

\section{When Recall Beats Relearning}
\label{sec:theory}

Consider a stream whose concepts recur in blocks of $L$ rows, batches of $B$ rows, labels one batch late, and a
context budget of $M$ rows. For a concept $k$ let $a_k(n)$ be the accuracy of the frozen TFM with a context of $n$ rows
of that concept on fresh rows of it (its \emph{learning curve}), non-decreasing in $n$.

\begin{assumption}\label{ass:id}
After the first labelled batch of a block, the meta expert puts its weight on the stored expert of the active
concept (identification takes one batch), and a reset learner (drift detection) starts its context at the block
boundary.
\end{assumption}

\begin{proposition}[Recall versus relearn]\label{prop:recall}
Under \cref{ass:id}, on every recurrence of concept $k$ after its first visit, the expected accuracy of recall minus
that of relearning, averaged over the $J=L/B$ batches of the block, is
\[
  \Delta_k(L)=\frac1J\sum_{j=1}^{J-1}\Bigl[a_k(M)-a_k\bigl(\min(jB,M)\bigr)\Bigr]\;\ge\;0,
\]
the area between the learning curve and its value at the budget. It vanishes when $a_k(B)=a_k(M)$, satisfies
$\Delta_k(L)\le\frac{M}{L}\bigl(a_k(M)-a_k(B)\bigr)$, so it decays as $O(M/L)$ for long blocks, and with $C$ cycles
over the concepts the overall gain is $\frac{C-1}{C}\Delta_k(L)$.
\end{proposition}

\begin{proof}[Proof sketch]
Both learners lose the first batch of the block, since its labels have not arrived. In batch $j+1$ of the block
($j=1,\dots,J-1$), the reset learner has the $\min(jB,M)$ labelled rows of the block in its context, and the recalling
learner has the stored context of $M$ rows of concept $k$ (\cref{ass:id}); taking expectations over fresh rows gives
the bracket. Monotonicity of $a_k$ gives $\Delta_k\ge0$. Every bracket is at most $a_k(M)-a_k(B)$ and vanishes for
$jB\ge M$, so at most $M/B$ of the $J=L/B$ terms are non-zero, which gives the $O(M/L)$ bound. The first visit has
no stored context, which gives the factor $(C-1)/C$.
\end{proof}

\Cref{prop:recall} makes the value of memory a property of the backbone and the concept, measurable before deployment
from a learning curve. The identification assumption is optimistic; the exponential-weights regret of the meta
expert, $O(\sqrt{T\ln K})$ \citep{zhao2025mooe}, bounds what it costs when the active expert is not identified at
once.\todo{Replace the assumption by the regret bound; prove the per-row monotonicity in $L$.}

% =====================================================================================================
\section{Experiments}
\label{sec:exp}

\textbf{RQ: does the learning curve predict the value of memory?} On a controlled grid of recurring concepts (RBF
mixtures with 5, 30 or 100 centroids whose class assignment is permuted between three concepts, so that $P(\vx)$ is
fixed; blocks of 500 or 2000 rows; three cycles), we measured the TFM's learning curves on held-out seeds, wrote down
the gain predicted by \cref{prop:recall}, and only then ran the streams once (\cref{tab:grid}). Memory beats a
reset-and-relearn detector on all eight streams with 30 or 100 centroids, more for shorter blocks and more complex
concepts, as predicted; the rank correlation between predicted and observed gain over the six cells is 0.94. A
window ensemble without memory does not reach \method{} on any of these streams, so the gain is due to recall rather
than faster adaptation. With 5 centroids the prediction is zero; one of four streams shows a gain of 5.3 points that
comes from a slow detection of the reset baseline, not from memory.

\begin{table}[t]
  \centering
  \caption{Recall versus relearn on the controlled grid (held-out seeds 10 and 11, mean). Gain = accuracy of \method{}
  minus that of a DDM-triggered reset; predicted = \cref{prop:recall} from learning curves measured before the run.}
  \label{tab:grid}
  \small
  \begin{tabular}{lcccc}
    \toprule
    Centroids & Block & Predicted & Gain & Oracle \\
    \midrule
    5   & 500  & 0.001 & 0.022 & 0.964 \\
    5   & 2000 & 0.000 & $-$0.000 & 0.992 \\
    30  & 500  & 0.019 & 0.023 & 0.937 \\
    30  & 2000 & 0.006 & 0.005 & 0.971 \\
    100 & 500  & 0.104 & 0.086 & 0.823 \\
    100 & 2000 & 0.030 & 0.026 & 0.897 \\
    \bottomrule
  \end{tabular}
\end{table}

\begin{itemize}
  \item \textbf{RQ1}: Do FIFO contexts forget recurring concepts (Observations~\ref{obs:forget},~\ref{obs:relearn})?
  \item \textbf{RQ2}: Does \method{} improve prequential accuracy, in particular right after a recurrence?
  \item \textbf{RQ3}: Do label-free prior weights recall the right expert before labels arrive?
  \item \textbf{RQ4}: Are new concepts detected (delay, false alarms)?
  \item \textbf{RQ5}: Cost: memory of the expert pool and time per instance.
  \item \textbf{RQ6}: Ablation of the prior weights, the meta expert and expert merging.
\end{itemize}
\textbf{Data.} Synthetic streams with recurring concepts (SEA, Agrawal, STAGGER, RBF with periodic switching); real
streams used by \citet{lourenco2026streams} (Electricity, NOAA weather, Forest CoverType, PokerHand, Posture, METER).
\textbf{Baselines.} TabPFN with a sliding window and with the dual memory of \citet{lourenco2026streams};
Drift-Resilient TabPFN \citep{helli2024drift}; MOOE with trained experts \citep{zhao2025mooe}; Adaptive Random Forest,
Streaming Random Patches and recurring-concept methods from stream mining.\todo{Pick recurring-concept baselines
(e.g.\ RCD, ensemble pools) after a literature check.} \textbf{Backbones.} TabPFN v2, TabICL.

\section{Related Work}
Concept drift \citep{gama2014survey,lu2019learning}. Recurring concepts have been handled by storing past models and
reusing them, selected by statistical tests \citep{goncalves2013rcd} or by a meta-learner of each model's domain of
applicability that can detect recurrence from unlabelled rows \citep{gama2014recurrent}; \method{} stores contexts
instead of trained models, so storing a concept costs no training, and the value of storing it is predicted by the
learning curve; concept drift for new model families,
e.g.\ multi-modal LLMs \citep{yang2025adapting,yang2026turning}; online and offline experts
\citep{zhao2025mooe}; TFMs on streams and under temporal shift \citep{lourenco2026streams,helli2024drift}.
\textbf{Difference to the competitors.} \citet{lourenco2026streams} keep one FIFO short- and long-term memory and do
not model recurrence; Drift-Resilient TabPFN retrains the prior to extrapolate temporal trends. \method{} stores past
contexts as experts, recalls them before labels arrive, and has regret guarantees.


\appendix
\section{Design History and Pre-Registration}
\label{app:history}
Development used synthetic streams with seeds 0--2 and the grid with seeds 0--1. The first version weighted experts
with $\eta=2$ and a discounted loss ($\gamma=0.5$) and had a single window; it improved on FIFO but not on a
reset-on-drift detector or on a window ensemble without memory, which reach the label-delay ceiling on simple concepts.
This led to the recall-versus-relearn question, to adding windows of 100 and 300 rows, and to choosing $\eta=10$ on
the last labelled batch on development streams. Expert predictions are cached, and the frozen weighting is replayed
causally (it only uses losses of batches whose labels have arrived), which equals running the method online. Learning
curves and predicted gains on the held-out grid seeds were written down before the grid was run; the grid hypotheses
held except an incidental expectation of no gain with 5 centroids, missed on one stream because of the baseline.

\bibliography{refs}
\bibliographystyle{icml2026}

\end{document}
